\chapter{Cross-spectra, coherence, and Granger prediction}\label{ch:connectivity}
\begin{readergoal}Derive dependence measures for multivariate signals and understand why prediction is not intervention.\end{readergoal}
\section{From autocovariance to cross-covariance}
For weakly stationary processes $X_t$ and $Y_t$, define $C_{xy}(\tau)=E[(X_t-\mu_x)(Y_{t+\tau}-\mu_y)]$. Its Fourier transform is the cross-spectral density $S_{xy}(f)$. The autospectra $S_{xx},S_{yy}$ describe each process; the cross-spectrum is complex and contains magnitude and phase information about their lagged relationship.

Magnitude-squared coherence is
\[\mathcal C_{xy}(f)=\frac{|S_{xy}(f)|^2}{S_{xx}(f)S_{yy}(f)},\]
when the denominator is positive. For valid spectral-density matrices, $0\le\mathcal C_{xy}(f)\le1$. Coherence is non-directional. Magnitude-only summaries discard phase and cannot by themselves indicate whether one observed signal drives another.

\section{Common input and reference contamination}
Suppose $x=s+n_x$ and $y=s+n_y$ share a signal $s$ while $n_x,n_y$ are independent noise processes. Then $x$ and $y$ can exhibit substantial coherence even though no direct interaction between their unique sources exists. Similarly, a reference electrode that carries electrical activity can insert a common component into many recordings. Bipolar referencing or spatial derivatives may reduce some effects but introduce others; one must evaluate the measurement geometry and signal-to-noise ratio \cite{Bastos2016,CohenTsuchiya2018}.

\section{Autoregressive models and direction of prediction}
A vector autoregressive model of order $p$ for $z_t\in\mathbb R^d$ is
\[z_t=\sum_{k=1}^p A_kz_{t-k}+\varepsilon_t.\]
Consider predicting $Y_t$ from its own past versus predicting it from both its own past and the past of $X_t$. A classical Gaussian linear Granger statistic compares residual variances:
\[F_{X\to Y}=\log\frac{\sigma^2_{\rm restricted}}{\sigma^2_{\rm full}}.\]
When the full model predicts better, the statistic is positive. The arrow denotes an improvement in prediction under the candidate model, not proof that manipulating $X$ would change $Y$. Latent common causes, instantaneous mixing, nonstationarity, model order, and noise can alter apparent direction. Spectral Granger extensions decompose predictive relationships by frequency under additional factorization assumptions \cite{Cekic2018}.

\section{Transfer entropy and conditional dependence}
Transfer entropy is a conditional information measure, often written $I(Y_t;X_{t-1}^{(p)}\mid Y_{t-1}^{(q)})$. It compares how informative the past of $X$ is about $Y_t$ once the past of $Y$ is known. Nonlinear estimators can be useful but require much data and careful treatment of bias. Prediction-based asymmetry is not a universal causal-identification guarantee.

\begin{workedexample}[A shared driver]
Let $s_t$ be a stochastic signal, $x_t=s_t+\eta_t$, and $y_t=s_{t-1}+\xi_t$ with independent observation noises. The lag creates directional predictability from $x$ to $y$, but the actual source is $s$, not a connection from the observed $x$ channel to the observed $y$ channel. This is a constructive counterexample to an unrestricted causal interpretation.
\end{workedexample}
\begin{exercises}
\item State why coherence is unchanged by a fixed nonzero scaling of one signal.
\item Distinguish phase locking, coherence, and Granger prediction.
\item Construct an example of common reference that raises coherence without coupling.
\item Describe checks for VAR residual autocorrelation and stability.
\end{exercises}